\section{Cycle packing with no positive-excess parameter}
\label{sec:four-unrestricted-iterative-packing}

Throughout this section, $\xi_G$ is a uniform member of the binary cycle space of $G$.
The rank of a graph is its cycle-space dimension.
All logarithms have base two.
We use only the following qualitative input from
Lemma~\ref{lem:path-qualitative-fixed-ratio-core}:
there is an absolute $D_*\ge3$ such that, with $H=256$, every connected normalized pole of rank at least $H$, girth at least $D\ge D_*$, and at most $4D$ degree-two ports has forest probability at most $3/5$ in every feasible affine parity class.
All its other vertices have degree three.

\begin{lemma}[Terminal bipartitions along a forest]
\label{lem:four-terminal-bipartition}
Let $Q$ be a connected loopless multigraph and let $A$ be a nonempty proper vertex set.
For each vertex $v$, let $L_v$ be the span of its incident coordinate forms on $\mathcal Z(Q)$.
Put $U=\sum_{a\in A}L_a$ and $B=|\delta_QA|$.
Among vertices $v\notin A$ with $L_v\cap U\ne0$, at most $2B$ are exceptional.
The others form at most $2B$ ordered classes.
Within each class, the vertices are articulation vertices of one component of $Q-A$ and give the same bipartition of its boundary terminals.
Each edge to $A$ supplies a distinct terminal.
\end{lemma}

\begin{proof}
Exclude the at most $B$ vertices adjacent to $A$.
In each component of $Q-A$, choose a spanning tree, attach a leaf for each boundary terminal, and retain the union of paths between those leaves.
This forest has at most $B$ leaves and at most $B$ branching vertices.
Exclude its branching vertices and suppress its degree-two vertices.
The resulting forest has at most $2B$ edges.

We verify that those edges contain all other relevant vertices.
Fix such a vertex $v$ and choose $0\ne a\in L_v\cap U$.
Represent $a$ both by a subset of the star at $v$ and by edges incident with $A$.
Their symmetric difference annihilates the cycle space, so it is a cut.
This elementary cut-space fact follows, for example, by assigning each vertex the parity of a root-to-vertex tree path; orthogonality to every fundamental cycle gives the same endpoint rule on each nontree edge.

Let $K$ be the component of $Q-A$ containing $v$.
Every component of $K-v$ lies entirely on one side of that cut.
Thus the star subset joins $v$ to a union of components of $K-v$.
As a functional, it equals the sum of the terminal coordinates in that union.
The sum of all terminal coordinates of $K$ is zero.
Since $a\ne0$, at least two components of $K-v$ contain terminals.
Every spanning tree must connect those terminal groups through $v$.
After removing the branching exceptions, exactly two terminal-bearing components remain.
Their terminal sets are the two sides of one edge of the suppressed forest.
Vertices on the same edge therefore give the same actual terminal bipartition.
Ordering them along that edge proves the assertion.
\end{proof}

\begin{lemma}[Two clean regions]
\label{lem:four-two-clean-regions}
Let $G$ be connected and cubic, let $P$ be connected or empty, and let $S$ be disjoint from $P$.
Suppose $G-(P\cup S)$ has girth at least $D\ge D_*$.
If $R_1,R_2$ are disjoint connected vertex sets avoiding $P\cup S$, each of rank at least $H$ and boundary at most $4D$, then
\[
 \Pr(\xi_G-(P\cup S)\text{ is a forest})\le\frac9{25}.
\]
\end{lemma}

\begin{proof}
Prune leaves from each induced graph $G[R_i]$.
Its connected two-core has the same rank and no larger boundary.
Indeed, removing a leaf in a cubic ambient graph replaces two external incidences by one, decreasing the boundary by one.
Each core has internal degrees two and three, with one boundary edge at each degree-two vertex.
It retains girth at least $D$.

Condition on every edge outside the two core interiors, including any edges between the cores.
For every conditioning of positive probability, the internal fiber is the Cartesian product of two affine translates of the internal cycle spaces.
Thus the two internal states are independent and uniform in their parity classes.
Odd prescribed parity can occur only at degree-two ports, and its port set has even size.
Each conditional forest probability is at most $3/5$ by the qualitative core lemma.
If a core has zero or one port, its only feasible parity class is the even space, whose only forest is empty and has probability $2^{-h}\le3/5$.
The global forest event requires both restrictions to be forests, so its probability is at most $(3/5)^2$.
\end{proof}

\begin{proposition}[Packing after a previous deletion]
\label{prop:four-packing-after-deletion}
Under the preceding notation, suppose
\[
 \Pr(\xi_G-(P\cup S)\text{ is a forest})>\frac12,
 \qquad s=|S|.
\]
Let $2\le d\le2D$.
Every family of $q$ disjoint connected cyclic regions in $G-(P\cup S)$, each with boundary in $G$ at most $d$, satisfies
\[
 q<(6H+15)d(s+1)2^d.
\]
\end{proposition}

\begin{proof}
If $q=0$ there is nothing to prove.
Contract each region to a marked vertex, and contract each connected component outside all regions to one unmarked vertex.
Discard loops, and call the resulting connected quotient $Q$.
Its marked vertices have degree at most $d$.

Restriction to the quotient edges maps $\mathcal Z(G)$ onto $\mathcal Z(Q)$.
To see surjectivity, extend the even quotient boundary pattern inside each connected contracted part by any parity join.
The fiber is a product of affine internal cycle spaces, so the quotient state is uniform and its internal choices are conditionally independent.
At a marked vertex $v$, let $I_v$ indicate that every incident quotient edge is absent.
In that event the internal state is uniform even; since the region is cyclic, being a forest requires the empty state and has probability at most $1/2$.
When $I_v=0$, use the upper bound one.
It follows that
\[
 \Pr(\xi_G-(P\cup S)\text{ is a forest})
 \le\mathbb E\prod_v(1-I_v/2).
\]

Write $L_v$ for the incident coordinate span on $\mathcal Z(Q)$.
It has dimension at most $d-1$, because even star parity gives one relation.
No $L_v$ is zero: a zero space would make $I_v=1$ identically and the displayed bound at most one half.
Take a maximal family of these spaces forming a direct sum, with vertex set $A$ and size $t$.
For this family the zero-star events are independent, each with probability
$2^{-\dim L_v}\ge2^{1-d}$.
Discarding the other factors gives
\[
 \Pr(\xi_G-(P\cup S)\text{ is a forest})
 \le(1-2^{-d})^t.
\]
If $t\ge2^d$, this is at most $e^{-1}<1/2$.
Consequently $1\le t<2^d$.
If all marked vertices are selected, the required conclusion follows already.

Otherwise put $U=\sum_{a\in A}L_a$ and $B=|\delta_QA|\le dt$.
Maximality gives $L_v\cap U\ne0$ at every remaining marked vertex.
Lemma~\ref{lem:four-terminal-bipartition} places all but at most $2B$ of them in at most $2B$ ordered classes.

Consider one class $v_1,\ldots,v_m$.
Every consecutive block of $H+2$ class vertices gives a middle component of $Q-A$ after deleting its first and last vertices.
It contains the $H$ interior marked vertices.
It has no neighbor in $A$: any attached terminal in that middle component would change the actual terminal bipartition between the first and last vertices.
Every edge leaving it therefore ends at those two vertices, giving boundary at most $2d$.
The articulation separations make the middle components from disjoint consecutive blocks pairwise disjoint.
Their lifts to $G$ are connected and contain $H$ disjoint cyclic regions, hence have rank at least $H$.

If $P$ is nonempty, its connectedness and avoidance of all marked regions put it in one contracted outside component.
At most one middle lift can meet it.
At most $s$ further lifts meet $S$.
Thus $m\ge(H+2)(s+3)$ would give at least two clean lifts, each with boundary at most $2d\le4D$.
Lemma~\ref{lem:four-two-clean-regions} would make the forest probability at most $9/25$, a contradiction.
Every class therefore has fewer than $(H+2)(s+3)$ members.

Counting selected vertices, exceptions, and classes now gives
\begin{align*}
 q&\le t+2B+2B(H+2)(s+3)\\
  &\le t+(6H+14)(s+1)B\\
  &\le(6H+15)d(s+1)t
   <(6H+15)d(s+1)2^d.
\end{align*}
\end{proof}

\begin{theorem}[An excess-free girth-raising bound]
\label{thm:four-excess-free-girth-raising}
There is an absolute constant $K$ with the following property.
Let $G$ be connected and cubic and let $P$ be connected or empty.
Assume
\[
 \Pr(\xi_G-P\text{ is a forest})>\frac12,
 \qquad \operatorname{girth}(G-P)\ge D_*.
\]
For each integer $D\ge D_*$, there is $S\subseteq V(G)\setminus P$ such that
\[
 \operatorname{girth}(G-(P\cup S))>D,
 \qquad |S|+1\le2^{\,2D+K\log^2(D+2)}.
\]
The set $S$ is the vertex union of a disjoint cycle family, maximal among cycles of length at most $D$ outside $P$.
On the common forward dyadic scales $D_j=2^jD_*$, the sets $S_j$ can be chosen nested with this bound simultaneously at every scale.
\end{theorem}

\begin{proof}
Deletion can only increase the probability of the forest event.
To raise an available girth scale $a\ge D_*$ to $b\le2a$, take a maximal disjoint family of cycles of length at most $b$ outside $P\cup S$ and add its vertices to $S$.
Every cycle has boundary at most its length in the cubic ambient graph.
The packing proposition, applied with $d=b$, gives
\[
 s_{\mathrm{new}}+1
 \le(6H+16)b^2 2^b(s+1).
\]
This also holds for an empty family.
The remaining graph has girth greater than $b$.

For an arbitrary target $D$, plan the scales backwards as $\lceil D/2^j\rceil$, stopping at the first value between $D_*$ and $2D_*$.
Read them forwards, starting with $S=\varnothing$ and available girth $D_*$.
If $D=D_*$, perform one deletion at that same scale.
Every step at most doubles the available scale; the sum of the new scales is at most $2D+O(\log(D+2))$.
There are $O(\log(D+2))$ steps.
Taking logarithms of the recurrence gives the asserted bound after increasing the absolute constant $K$.

Alternatively, use the common forward sequence $D_*,2D_*,4D_*,\ldots$.
The scales through $D_j$ sum to $2D_j-D_*$, so the same estimate holds simultaneously, and deletions are nested.
At every stage new cycles avoid all earlier deleted vertices.
Thus the combined family is disjoint, and the final girth assertion is exactly its maximality at the target length.
\end{proof}

\paragraph{Verification.}
This fragment restates the independently audited packing argument without its earlier forbidden-cut formulation.
An exact checker verifies the terminal partitions, middle intervals, core pruning, and complete conditional affine products:
\begin{center}
\small\nolinkurl{four-high-connectivity-iterative-packing-check.py}.
\end{center}
It does not numerically test the absolute high-girth threshold.
